\subsection{Efecto del problema en el área del trabajo o en los resultados en el trabajo}

El problema identificado en el proceso de gestión de nóminas y procesos administrativos de recursos humanos tiene un impacto directo en estas áreas de trabajo. 
La dependencia completa a las hojas de cálculo y la falta de un control centralizado de la información generan discrepancias entre registros, 
omisiones de incidencias y necesidad de corregir y recalcular constantemente, 
lo que afecta la eficiencia y la precisión de los resultados.

También se observa que la falta de seguridad y control de acceso expone la información sensible a posibles manipulaciones o pérdidas, 
lo que puede derivar en sanciones legales y pérdida de confianza por parte de los colaboradores. Con el propósito de mantener la trazabilidad metodológica del diagnóstico, en la Tabla~\ref{tab:causas_nominas} se codifican del CA1 al CA8 las ocho causas identificadas previamente en la problemática local (literales \textit{a} a \textit{h} del apartado~2.1.1), conservando su denominación y alcance para su posterior análisis de causa raíz y priorización.

% Tabla N1
\par\medskip
\noindent
\begin{minipage}{\linewidth}

    \refstepcounter{table}
    \label{tab:causas_nominas}
    \addcontentsline{lot}{table}{%
        Tabla~\thetable. Lista de causas identificadas}

    {\raggedright
        {\bfseries Tabla~\thetable}\par
        \vspace{0.1cm}
        {\itshape Lista de causas identificadas}\par
    }

    \vspace{0.2cm}

    \begingroup
    % Interlineado sencillo dentro de la tabla
    \renewcommand{\baselinestretch}{1}
    \small
    \renewcommand{\arraystretch}{1.25}

    \begin{tabularx}{\linewidth}{
        @{}p{0.12\linewidth}
        >{\raggedright\arraybackslash}X@{}
    }
        \toprule
        \textbf{Código} &
        \textbf{Identificación de las causas} \\
        \midrule

        CA1 &
        \textit{Almacenamiento de información laboral:} Información laboral y remunerativa almacenada en
        archivos separados entre Recursos Humanos y Contabilidad. \\

        CA2 &
        \textit{Control de versiones de archivos:} Circulación de distintas versiones de los archivos
        de planilla sin un control definido de actualización. \\

        CA3 &
        \textit{Transcripción de información:} Transcripción manual de asistencias, permisos, horas
        extras y otras incidencias laborales a la planilla. \\

        CA4 &
        \textit{Modificación de fórmulas de cálculo:} Modificación manual de fórmulas y parámetros de cálculo
        sin una validación uniforme. \\

        CA5 &
        \textit{Roles con permisos de usuario:} Permisos de consulta y modificación de datos laborales
        y salariales insuficientemente diferenciados según
        las responsabilidades del usuario. \\

        CA6 &
        \textit{Registro de auditoría:} Registro insuficiente de los cambios realizados sobre
        los datos laborales y los resultados de la planilla. \\

        CA7 &
        \textit{Preparación de pagos y declaraciones:} Preparación separada de información para pagos bancarios
        y declaraciones, mediante transferencias o registros manuales. \\

        CA8 &
        \textit{Generación y distribución de documentos:} Generación y distribución individual de boletas mediante
        actividades manuales o semimanuales. \\

        \bottomrule
    \end{tabularx}
    \endgroup

    \vspace{0.2cm}

    {\small\raggedright
        \textit{Nota.} Codificación de las causas identificadas en la problemática local (apartado 2.1.1).
        Elaboración propia.\par
    }

\end{minipage}
\par\medskip